\section{Introduction}
\label{sec:introduction}

Low-frequency acoustic and vibrational exposures are ubiquitous in
occupational and environmental settings---from wind turbines and heavy
machinery to transport vehicles and industrial plant.  The human body is known
to exhibit resonance phenomena in the \SIrange{4}{8}{\hertz} range under
whole-body vibration, with the abdominal viscera identified as the primary
resonant subsystem \cite{iso2631,griffin1990handbook}.

A persistent claim in popular culture holds that a specific infrasonic
frequency in the \SIrange{5}{9}{\hertz} range---the so-called ``brown
note''---can cause involuntary loss of bowel control through resonance of the
abdominal cavity.  Despite widespread repetition, no rigorous scientific study
has confirmed this claim, and it has been explicitly debunked in controlled
experiments \cite{Leventhall2007}.

However, the complete absence of a \emph{quantitative mechanistic analysis}
leaves the question incompletely resolved.  No study has:
\begin{enumerate}
  \item derived abdominal cavity eigenfrequencies from first-principles
    shell theory with proper fluid--structure coupling;
  \item quantified the airborne acoustic coupling coefficient for flexural
    modes in the long-wavelength limit; or
  \item compared airborne and mechanical excitation pathways using a common
    modal framework.
\end{enumerate}

The present work addresses all three gaps. We model the abdomen as a
fluid-filled viscoelastic oblate spheroidal shell and derive natural
frequencies for both breathing ($n=0$) and flexural ($n \geq 2$) mode
families.  We then compare the efficiency of airborne acoustic coupling
(subject to a $(ka)^n$ penalty) against mechanical whole-body vibration
coupling (which bypasses this penalty entirely).

The central finding is that the coupling disparity between these two
pathways---a factor of $10^3$--$10^4$---quantitatively explains why
whole-body vibration at \SIrange{4}{8}{\hertz} produces gastrointestinal
effects while airborne infrasound at comparable frequencies does not.

The remainder of this paper is organised as follows. Section~\ref{sec:formulation}
develops the mathematical formulation for the fluid-filled viscoelastic shell,
including both breathing and flexural mode families and the two coupling
pathways. Section~\ref{sec:results} presents a comprehensive parametric study
and validates the model against ISO~2631 data. Section~\ref{sec:coupling}
quantifies the airborne--mechanical coupling disparity.
Section~\ref{sec:discussion} interprets the results and discusses limitations,
and Section~\ref{sec:conclusion} summarises the conclusions.
